% vim: sw=0
\documentclass{resume}

\usepackage{hyperref}

\name{Minseo Kim (김민서)}

\email{\href{mailto:minseo@mskim.org}{minseo@mskim.org}}
\github{\href{https://github.com/kimminss0}{kimminss0}}
\blog{\href{https://blog.mskim.org}{blog.mskim.org}}
%\mobile{\href{tel:+821012345678}{010-1234-5678}}
\begin{document}

\begin{rSection}{Education}
	\begin{rEducation}{{\bfseries{POSTECH (포항공과대학교)}}}{Feb 2022 -- Present}
		B.S. in Computer Science and Engineering (GPA: 3.70/4.3)
	\end{rEducation}
\end{rSection}

\begin{rSection}{Skills}
	\begin{tabular}{@{}p{0.25\linewidth}p{0.7\linewidth}}
		Programming Languages
		 & C, C++, Python, Go, Java, TypeScript \\ [0.3em]

		System
		 & Linux, GDB                           \\ [0.3em]

		Infrastructure
		 & Docker, Docker Compose, Nginx        \\ [0.3em]

		Web
		 & React.js, Next.js, Express.js        \\ [0.3em]

		Databases
		 & PostgreSQL                           \\ [0.3em]

		Natural Languages
		 & Korean (Native), English             \\ [0.3em]
	\end{tabular}
\end{rSection}

\begin{rSection}{Job Experience}
	\begin{rSubsection}{\href{https://wordbricks.ai}{Wordbricks LLC}}{Jun 2025 -- Aug 2025}

		{\bfseries{Intern Full-stack Engineer}}

    \item LLM 기반 웹 크롤링 API 및 대시보드 개발 주도 (Product Owner)

		\item 기술 스택: TypeScript, Next.js, Hono, Vercel AI SDK, Cloudflare Workers, Prisma, Supabase
	\end{rSubsection}

	\begin{rSubsection}{\href{https://dslab.postech.ac.kr/}{POSTECH Data Systems Lab}}{Jan 2023 -- Jan 2024}

		{\bfseries{Undergraduate Researcher (Full-stack Web Developer)}}
    \item 국제 학술지 논문 검색 및 통계 웹사이트 개발

		\item 주요 기여 사항:

		\vspace{-0.5em}
		\begin{list}{--}{}
			\itemsep -0.5em
			\item 백엔드 및 프론트엔드 개발
			\item 데이터베이스 설계 및 운영
			\item 데이터 통합 및 파이프라인 설계
			\item 웹사이트 배포 및 서버 운영
		\end{list}

		\item 기술 스택: JavaScript, React.js, Express.js, PostgreSQL, Airflow, Docker, Docker Compose
	\end{rSubsection}
\end{rSection}

\begin{rSection}{Projects}
	\begin{rSubsection}{\href{https://github.com/kimminss0/CSED312} {Pintos}}{Sep 2025 -- Dec 2025}

		교육용 80x86 아키텍처 운영체제 개발 (Stanford CS140 기반)

		\item 주요 구현 사항:
		\vspace{-0.5em}
		\begin{list}{--}{}
			\itemsep -0.5em

			\item Kernel-level threading with priority scheduling and synchronization primitives (locks, semaphores)

			\item System call handling for user program execution, including process management and file operations

			\item Virtual memory system featuring demand paging, stack growth, memory mapped files
		\end{list}

		\item 사용 기술: C 및 x86 어셈블리

	\end{rSubsection}

	\begin{rSubsection}{\href{https://github.com/kimminss0/CSED311} {RISC-V CPU}}{Mar 2025 -- May 2025}

		5단계 파이프라인 RISC-V CPU (2인 프로젝트)

		\item RV32I 명령어 집합의 서브셋 구현 
		\item 주요 구현 사항:
		\vspace{-0.5em}
		\begin{list}{--}{}
			\itemsep -0.5em

			\item Hazard detection unit with data forwarding

			\item 2-Bit saturation counter branch predictor with pattern history table and globally shared history buffer (gshare)

			\item Direct-mapped cache
		\end{list}
		\item 사용 기술: Verilog HDL

	\end{rSubsection}

	\begin{rSubsection}{\href{https://github.com/kimminss0/CSED321/tree/main/hw7} {Tiny ML (TML) Compiler}}{May 2025}

    Standard ML (SML) 기반의 함수형 프로그래밍 언어 TML의 컴파일러 백엔드 구현

		\item 클로저(Closure), 고차 함수 및 재귀 함수, 재귀적 데이터 타입 및 패턴 매칭 지원
		\item 사용 기술: OCaml
	\end{rSubsection}

	\begin{rSubsection}{\href{https://github.com/kimminss0/CSED353} {Sponge}}{Mar 2025 -- May 2025}

		교육용 네트워크 스택 개발 프로젝트 (Stanford CS144 기반)

		\item 주요 구현 사항:
		\vspace{-0.5em}
		\begin{list}{--}{}
			\itemsep -0.5em

			\item TCP flow control

			\item Routing table lookup \& packet forwarding

			\item ARP \& IP packet processing
		\end{list}

		\item 사용 기술: C++
	\end{rSubsection}

	\begin{rSubsection}{\href{https://github.com/UniD3-Hackathon-Team4/barokey} {BaroKey}}{Nov 2023}

		제3회 유니디톤(UniDThon) 해커톤 출품작 (팀 프로젝트)

		\item 사용자 위치를 기반으로 안전 위협이나 자연재해 등 실시간 긴급 키워드를 제공하는 웹 서비스 개발

		\item 역할: {\bfseries{팀장}}, 풀스택 개발자 (프론트엔드 개발 기여, 백엔드 및 인프라 구축 리드)
	\end{rSubsection}
\end{rSection}

\end{document}
